%! Absolute Value Functions & Transformations — Unit 1, Lesson 1.3 — Homework (Answer Key)
% Mirrors homework/main.tex exactly; every \blank{} becomes an \ans{}, every \writelines{n} is
% answered with exactly n \ansline{}s, and every work block is authored byte-identically.
% No teacher notes here — they live in the lesson plan.
\documentclass[10pt]{article}
\usepackage{algebra2-article}
\usepackage{algebra2-key}   % pulls in -boxes; do NOT also load -boxes
% Coordinate grid: {xmin}{xmax}{ymin}{ymax}
\newcommand{\lgrid}[4]{%
  \draw[step=1, linegray!40, very thin] (#1,#3) grid (#2,#4);
  \draw[->, semithick, charcoal] (#1,0)--(#2,0) node[below right, font=\scriptsize]{$x$};
  \draw[->, semithick, charcoal] (0,#3)--(0,#4) node[left, font=\scriptsize]{$y$};}

\begin{document}

\pageheader{Unit 1, Lesson 1.3}{Homework --- Answer Key}
% NAMESTRIP: no \namedateperiod here — mirrors the blank.

\begin{remindbox}
\small
\textbf{This is your graded homework.} Your packet with completed homework is due the first
class after two study halls. I will announce the due date in class and on TurtleNet.
\end{remindbox}

\begin{notesbox}{Practice}
Show your work. In $a|x-h|+k$: vertex $(h,k)$ --- \textbf{set the inside to zero} --- axis $x=h$;
$a$ says up or down, $|a|$ says narrow or wide.

\begin{enumerate}[leftmargin=1.9em, itemsep=3pt]
  \item \textbf{Read the rule.} Fill in every feature.
    \begin{center}
    \renewcommand{\arraystretch}{1.3}
    \begin{tabularx}{\linewidth}{@{}l X X X X X@{}}
    & \textbf{Vertex} & \textbf{Axis} & \textbf{Opens} & \textbf{Max/min value} & \textbf{Range} \\
    (a) $y=|x-4|+2$ & \ans{$(4,2)$} & \ans{$x=4$} & \ans{up} & \ans{min $2$} & \ans{$y\ge2$} \\
    (b) $y=-\tfrac12|x+1|+3$ & \ans{$(-1,3)$} & \ans{$x=-1$} & \ans{down} & \ans{max $3$} & \ans{$y\le3$} \\
    (c) $y=3|x|-4$ & \ans{$(0,-4)$} & \ans{$x=0$} & \ans{up} & \ans{min $-4$} & \ans{$y\ge-4$} \\
    \end{tabularx}
    \end{center}

  \boxguard[18]
  \item \textbf{One sign apart.} Compare the two rules, and notice what changes.
    \begin{center}
    \renewcommand{\arraystretch}{1.3}
    \begin{tabularx}{\linewidth}{@{}X|X@{}}
    \textbf{(a) $y=3|x-1|-2$} & \textbf{(b) $y=-3|x-1|-2$} \\
    \hline
    Vertex \ans{$(1,-2)$}; opens \ans{up} & Vertex \ans{$(1,-2)$}; opens \ans{down} \\
    $-2$ is a \ans{minimum}; range \ans{$y\ge-2$} & $-2$ is a \ans{maximum}; range \ans{$y\le-2$} \\
    \end{tabularx}
    \end{center}
    (c) The two graphs share a vertex. Explain why that same point is a minimum for one and a
    maximum for the other.
    \par\ansline{With $a=3>0$ both rays rise away from the vertex, so every other output is above $-2$.}
    \par\ansline{With $a=-3<0$ both rays fall away from it, so every other output is below $-2$.}

  \boxguard[20]
  \item \textbf{Match, then complete a table.}
    \begin{enumerate}[label=(\alph*), leftmargin=1.9em, itemsep=4pt]
      \item Write the letter of each rule's graph.
        \begin{center}
        \begin{tikzpicture}[scale=0.30]
          \begin{scope}
            \lgrid{-1}{5}{-1}{4}
            \draw[very thick, forest] (-1,3)--(2,0)--(5,3);
            \node[charcoal, font=\scriptsize] at (2,-2.4){\textbf{W}};
          \end{scope}
          \begin{scope}[xshift=8.5cm]
            \lgrid{-3}{3}{-3}{3}
            \draw[very thick, forest] (-3,1)--(0,-2)--(3,1);
            \node[charcoal, font=\scriptsize] at (0,-4.6){\textbf{X}};
          \end{scope}
          \begin{scope}[xshift=17cm]
            \lgrid{-3}{3}{-4}{2}
            \draw[very thick, forest] (-3,-3)--(0,0)--(3,-3);
            \node[charcoal, font=\scriptsize] at (0,-5.6){\textbf{Y}};
          \end{scope}
        \end{tikzpicture}
        \end{center}
        $y=|x-2|$ is \ans{W} \quad $y=|x|-2$ is \ans{X} \quad $y=-|x|$ is \ans{Y}
      \item Complete the table for $y=|x-1|+2$, then read the vertex off it.
        \begin{center}
        \renewcommand{\arraystretch}{1.2}
        \begin{tabular}{|c|c|c|c|c|c|}
        \hline
        $x$ & $-1$ & $0$ & $1$ & $2$ & $3$ \\ \hline
        $y$ & \ans{$4$} & \ans{$3$} & \ans{$2$} & \ans{$3$} & \ans{$4$} \\ \hline
        \end{tabular}
        \end{center}
        Vertex \ans{$(1,2)$} --- the row where $y$ stops falling and starts rising.
    \end{enumerate}

  \boxguard[16]
  \item \textbf{The odd ones out.} Answer each from the picture, not a memorised rule.
    \begin{enumerate}[label=(\alph*), leftmargin=1.9em, itemsep=4pt]
      \item $y=|x|+k$ crosses the $x$-axis \ans{two} times when $k<0$, \ans{one} time(s)
            when $k=0$, and \ans{zero} times when $k>0$.
      \item For which values of $a$ is $y=a|x|$ \emph{wider} than the parent? \ans{$-1<a<1$, $a\ne0$}
            \; What happens to the graph when $a=0$? \ans{it flattens to the line $y=0$ --- no V at all}
    \end{enumerate}

  \boxguard[22]
  \item \textbf{Model it.} A drone flies in a straight line past a tower. Its distance from the
        tower, in metres, after $t$ seconds is $D(t)=15\,|t-6|$.
    \begin{enumerate}[label=(\alph*), leftmargin=1.9em, itemsep=4pt]
      \item The vertex is \ans{$(6,0)$}. In the story it means: \ans{at $6$ s the drone passes the tower (distance $0$)}
      \item When is the drone $45$ m from the tower? Solve $D(t)=45$.
        \begin{work}
          15\,|t-6| &= 45 \\
          |t-6| &= 3 \\
          t - 6 &= 3 \quad\text{or}\quad t - 6 = -3 \\
          t &= 9 \quad\text{or}\quad t = 3
        \end{work}
      \item If the drone flew \emph{twice as fast}, which number in the rule changes, and how does
            the V change?
        \par\ansline{$a$ doubles from $15$ to $30$ --- the distance grows twice as fast each second.}
        \par\ansline{The V gets narrower (steeper rays); the vertex $(6,0)$ does not move.}
    \end{enumerate}

  \boxguard[16]
  \item \textbf{Multiple choice (SOL-style).} Which describes how the graph of $y=|x+4|-2$ is
        obtained from the graph of $y=|x|$?
    \par\vspace{2pt}\noindent
    \begin{tabular}{@{}p{0.46\linewidth}p{0.46\linewidth}@{}}
      A.\; Shift right $4$ and down $2$ & B.\; \ans{Shift left $4$ and down $2$} \\
      C.\; Shift left $4$ and up $2$    & D.\; Shift right $4$ and up $2$ \\
    \end{tabular}
    \par\vspace{2pt}\noindent
    Say in one sentence how you ruled out one of the others.
    \par\ansline{B. Inside $=0$ at $x=-4$: the vertex is at $x=-4$, a shift \emph{left} --- so A and D are out.}
    \par\ansline{The $-2$ is outside the bars, so it moves the vertex \emph{down} --- so C is out.}
\end{enumerate}
\end{notesbox}

\boxguard[18]
\begin{extensionbox}
\textbf{Build and reason.}
\begin{enumerate}[label=(\alph*), leftmargin=1.9em, itemsep=2pt]
  \item A V opens \emph{up}, has vertex $(-2,1)$, and passes through $(0,5)$. Find $a$ and write
        the rule in vertex form.
        \begin{work}
          5 &= a\,|0-(-2)| + 1 \\
          5 &= 2a + 1 \\
          a &= 2 \\
          y &= 2|x+2| + 1
        \end{work}
  \item A classmate claims that $y=-2|x-3|+7$ reaches the output $9$ somewhere. Explain from the
        picture why that cannot happen.
        \par\ansline{$a=-2<0$: the V opens down, so the vertex $(3,7)$ is its highest point; range $y\le7$.}
        \par\ansline{Both rays fall away from $(3,7)$, so no output is above $7$ --- $9$ is unreachable.}
\end{enumerate}
\end{extensionbox}

\begin{spiralbox}
\textbf{Coming next --- Lesson 1.4: Piecewise \& Step Functions.} Today's V is secretly
\emph{two} lines pasted at the vertex. Next you write functions built from \emph{different rules
on different pieces of the domain} --- the V is the first one you know --- and meet the step
function, whose graph is a staircase.
\end{spiralbox}

\end{document}
